\section{Motivation}

Multi-agent scenarios are ubiquitous in everyday life. We often need to make decisions based on information on the problem to be solved and on our knowledge, or belief, about the preferences and actions of other agents involved in, or impacted by, the decision.
It is therefore an epistemic reasoning task to make decisions in this scenarios.

Often such decision involve creating a plan, that is, a sequence of actions that, when executed, will lead to a desired goal state.
The research area of \emph{Multi-agent Epistemic Planning} (MEP) attempts to build artificial agents that can build such plans. 
%has recently gained attention within the planning community.
MEP is a generalization of the concept of \emph{Epistemic Planning} (EP),
which refers to a single agent setting that reasons about knowledge and beliefs. In the MEP scenario, agents also reason about 
%, MEP refers to a setting of automated planning  where  agents reason in the space of knowledge states \Lior{I know that this comes later, but it may be useful to either contrast with non-epistemic planning, with respect to the space they operate in (not a knowledge space), or provide a forward reference} and on 
the information that other agents have of each other, and how such information flows between agents.
%, and try to reach a desirable state from the current state. 
% More and more researchers are approaching such formalism that allows to represent scenarios in which the agents need to reason on both the state of the world and on the information flow.

%he ability of representing knowledge or belief of agents and to plan while considering (and manipulating) them has, in fact, captured the interest of many researchers.
%This interest has led to improvements in theoretical and practical aspects of multi-agent epistemic planning.
%In particular, s
Several approaches to develop epistemic solvers/planners have been made across the community, see for example~\cite{baral2015action,muise2015planning,Fab20}.
While these solvers may have different restrictions, \eg considering only the concept of \emph{knowledge} while not allowing the idea of \emph{belief}, or not allowing for \emph{dynamic common \knobel}, all of them share the same goal: to plan while considering the information flows between the agents.

While the diversification of approaches in solving this problem certainly brings benefits to the entire community, with many specialized efficient planners that are suitable for specific scenarios, the lack of a standardized way to define input problems in multi-agent epistemic planning has led to the creation of ad-hoc languages that can become an obstacle 
%\Lior{can become an obstacle/ confer an obstacle} 
when comparing, combining (i.e., ensemble), or testing different planners.
We believe that a unified way of expressing MEP problems would provide researchers with a faster and less error-prone way of defining standardized benchmarks and would also relieve them from the burden of having to learn a new language each time they need to make use of a different tool.
Moreover, it would allow for easy transfer of techniques and results between different planning scenarios and planner methodologies.


In this paper we therefore propose the Epistemic Planning Domain Definition Language (E-PDDL), a unified way of defining epistemic problems that, similarly to the well-known language PDDL for classical planning, could be regarded as the standard 
%action 
language to be used to express MEP problems.
Being able to define MEP problems in E-PPDL and then to map them into 
specific languages for existing planners provides also the necessary 
environment for a real-time choice of the best planner for a given MEP problem. Indeed, this work is part of a larger effort that aims to equip artificial agents with meta-cognitive capabilities that, given internal values, resources, and criteria, can identify some of the problem features and assess what skills are required to best solve it. In the context of this paper, this would mean to be able to select the most suitable planner for a MEP problem. This is greatly facilitated if there is a way to map the input problem specification into the language required by each specific planner.

In the rest of the paper, we  first give necessary background about EP and MEP with a running example. We then introduce the E-PDDL language and its key features, and we show how it encompasses current MEP formalisms. Next, we highlight implementation details in the parser and conclude.
